\documentclass[conference]{IEEEtran}

\usepackage[T1]{fontenc}
\usepackage{amsmath,amssymb}
\usepackage{booktabs}
\usepackage{graphicx}
\usepackage{cite}
\usepackage{tikz}
\usetikzlibrary{arrows.meta,positioning}
\usepackage[hidelinks]{hyperref}

% This file is the report structure. Replace TODOs with reviewed content.
\newcommand{\placeholder}[1]{\noindent\textit{[TODO: #1]}}

\begin{document}
\title{Verifier-Guided Revision for Instruction Following: Comparing Feedback Granularity on IFBench}
\author{\IEEEauthorblockN{[Author name]}
\IEEEauthorblockA{COMP4136 Natural Language Processing\\Hong Kong Baptist University}}
\maketitle

\begin{abstract}
Large language models can answer a question while missing or disobeying its output constraints, and revision can also introduce new violations. We evaluated whether verifier feedback supported more reliable revision on IFBench. Using two models, gpt-4.1 and gpt-4o, we compared three revision branches starting from a shared initial answer (B0): unaided self-revision (B1), revision with overall verifier feedback (B2), and revision with per-constraint feedback (B3). Each branch started independently from the same initial answer. Compared with B0, B2 increased paired strict prompt-level accuracy by 9.65 and 10.04 percentage points for the two models respectively. Both differences remained statistically significant in post-hoc cluster sign-flip tests with Holm correction. B1 had higher observed regression rates among initially passing answers than B2, while the comparison between B2 and B3 did not show a statistically significant difference. A content review covered all 260 test prompts and found that B2 achieved completion rates of 99.2\% for gpt-4.1 and 98.5\% for gpt-4o. After post-hoc paired comparisons, no statistically significant differences in completion rates were recognized between the revision conditions and B0 after Holm correction. Overall, verifier feedback improved strict instruction following in this experiment.
\end{abstract}

\section{Introduction}
Instruction following requires a model to complete the requested task as well as satisfying the user's output constraints. For example, a mathematical explanation may be correct but exceed a word limit, while a response that follows a lexical rule may still provide an inadequate explanation. Therefore, these properties should be assessed separately. IFEval introduced executable checks for verifiable instructions~\cite{zhou2023ifeval}, and IFBench extended this approach to constraint types that models are unfamiliar with~\cite{pyatkin2025ifbench}. These checks can help inspect individual constraint failures, and their results can be used as feedback for model's further revision and generation.

Revision provides models with a second chance to satisfy an output constraint, but it can also damage initially passing answers. Thus, all cases where the original answer was correct or not need to be considered for the experiment. Our experiment simultaneously maintained both types of responses and separately measured transitions from failure to success as well as from success to failure.

The central question this project addressed was whether overall verifier feedback and per-constraint feedback led to different revision outcomes. We compared three independent revision branches against a shared initial answer: self-revision without external feedback (B1), revision with an overall strict pass/fail result (B2), and revision with a strict result for each constraint (B3). We used two models named gpt-4.1 and gpt-4o to examine whether the observed patterns appeared in both models.

Our first contribution is a paired implementation that preserved the original prompts, branch independence, verifier outcomes, and recoverable generation records. Our second contribution is an analysis that separates net accuracy changes from repairs and regressions. Our third contribution is a content review connecting rule outcomes to human-confirmed completion labels. Together, these contributions address whether feedback helps, how it helps, and what the verifier score leaves unmeasured on the instruction following scenario.

The results suggest that overall verifier feedback is a practical method to improving instruction following in this experiment. B2 increased paired strict accuracy by 9.65 percentage points for gpt-4.1 and 10.04 percentage points for gpt-4o. Its observed regression rates were only 5.1\% and 6.8\%, compared with 13.8\% and 10.8\% for B1. However, the evidence does not establish that per-constraint feedback was inferior because the B2 and B3 comparisons did not show statistically significant differences.

\section{Background Survey and Related Work}
\subsection{Evaluating Instruction Following}
A response can address the user's question and still violate a specified output constraint. For example, an answer may provide a relevant explanation while exceeding a word limit. To consistently evaluate such requirements, IFEval introduced prompts with instructions that can be checked by programs, including word-count and keyword requirements~\cite{zhou2023ifeval}. An automatic verifier makes the rule decision reproducible for a given response and implementation. However, it does not establish whether the response is accurate, sufficiently detailed, or useful.

IFBench extended the evaluation of verifiable instruction following to 58 out-of-domain constraint types~\cite{pyatkin2025ifbench}. Its focus on unfamiliar constraints makes it suitable for examining whether models can follow requirements beyond a small set of familiar benchmark rules. We used the IFBench prompts and its official strict and loose evaluators with the dataset version containing 300 prompts.
\subsection{Revising Answers with Feedback}
A model can revise its initial answer based on feedbacks generated by itself. Self-Refine uses this approach to refine answers over successive iterations~\cite{madaan2023selfrefine}. Our B1 branch is a simpler comparison version of Self-Refine which only asks the model to check its initial answer once and return a revision.

Feedback can also come from outside the model. CRITIC uses external tools to check aspects of an initial output and then uses their feedback for revision~\cite{gou2024critic}. In our experiment, we learned the revision logic from CRITIC and set IFBench verifier as the fixed evaluating tool to provide models feedbacks for revision.

\section{Methodology}
The revision workflow is shown in Fig.~\ref{fig:method}.
\subsection{Shared Initial Answers and Independent Revisions}
We used a shared initial answer (B0) to control the same starting point for the three revision branches. For each test prompt and model, B0 was generated from the original IFBench user message without an additional answer-format wrapper. Then, we constructed B1, B2, and B3 as independent conversation branches, containing the original message, the B0 assistant answer, and a revision instruction to produce one revised answer.

We revised both initially passing and failing answers because they were all relevant to the research question and meaningful to future analysis. We also had stored prompt identities, model identifiers, conditions, generation status, responses, and scores separately to preserve the relationship between each response and its experimental condition and supported the paired analysis.
\subsection{Feedback Conditions}
\textbf{B1: Self-Revision.} The model received the instruction: ``Check the original output constraints yourself. Do not assume they were already satisfied.'' This condition served as a control to observe the effect via providing the model with an additional generation opportunity under the general inspection request would have.

\textbf{B2: Revision with Overall Verifier Feedback.} B2 received the overall strict verifier outcome. A failed answer received: ``Verifier (strict): FAILED. At least one output constraint was not satisfied.'' A passing answer received: ``Verifier (strict): PASSED all output constraints.''

\textbf{B3: Revision with Per-Constraint Feedback.} B3 received the strict outcome for each constraint in its original order. Each line contained PASS or FAIL followed by the corresponding requirement extracted from the prompt.
\subsection{Prompt Design}
All revision branches were instructed to review the previous answer against the original user message, preserve the original task and its constraints, and return only the final answer. This design aimed to avoid explanatory headings or additional formatting that could violate the original constraints.

When the verifier reported that B0 had passed, B2 and B3 also allowed the model to leave the answer unchanged. The full prompt templates and message order are provided in Appendix~\ref{app:prompts}.
\begin{figure*}[!t]
  \centering
  % Vector diagram for the independent IFBench revision branches.
\definecolor{bzero}{HTML}{0072B2}
\definecolor{bone}{HTML}{E69F00}
\definecolor{btwo}{HTML}{009E73}
\definecolor{bthree}{HTML}{CC79A7}
\definecolor{flowgray}{HTML}{45535E}
\begin{tikzpicture}[
    font=\sffamily\small,
    block/.style={draw=flowgray, line width=0.65pt, rounded corners=2pt,
      fill=white, align=center, minimum height=0.94cm, text width=2.15cm,
      inner sep=4pt},
    branch/.style={block, text width=2.9cm, minimum height=0.98cm},
    flow/.style={-{Stealth[length=2.2mm]}, draw=flowgray, line width=0.7pt},
    note/.style={font=\sffamily\scriptsize, text=flowgray, align=center}
  ]
  \node[block] (prompt) at (0,0) {\textbf{IFBench}\\Original prompt};
  \node[block, draw=bzero, line width=0.9pt] (b0) at (3.1,0)
    {\textbf{B0}\\Initial answer};
  \node[block] (verify) at (6.2,-0.85) {\textbf{Strict}\\verifier};
  \node[branch, draw=bone, line width=0.9pt] (b1) at (10.0,1.5)
    {\textbf{B1} \; Self-revision\\Self-check request};
  \node[branch, draw=btwo, line width=0.9pt] (b2) at (10.0,0)
    {\textbf{B2} \; Overall feedback\\Pass / fail};
  \node[branch, draw=bthree, line width=0.9pt] (b3) at (10.0,-1.5)
    {\textbf{B3} \; Per-constraint feedback\\Pass / fail by constraint};
  \node[block, text width=2.35cm] (eval) at (14.0,0)
    {\textbf{Evaluation}\\Strict / loose scores};
  \draw[flow] (prompt) -- (b0);
  \draw[flow] (b0.north) |- (b1.west);
  \draw[flow] (b0.south) |- (verify.west);
  \draw[flow] (verify.north) |- (b2.west);
  \draw[flow] (verify.south) |- (b3.west);
  \draw[flow] (b1.east) -| (eval.north);
  \draw[flow] (b2.east) -- (eval.west);
  \draw[flow] (b3.east) -| (eval.south);
  \node[note] at (7.9,2.42) {All branches receive the same original prompt and B0 answer};
  \node[note] at (14.0,-2.25) {Content review\\conducted separately};
\end{tikzpicture}

  \caption{Experimental workflow. Each revision branch independently receives the original prompt and the same B0 answer. The strict verifier supplies overall feedback to B2 and per-constraint feedback to B3. All available responses are evaluated; content completion is reviewed separately.}
  \label{fig:method}
\end{figure*}

\section{Experimental Setup}
\subsection{Data and Protocol Freeze}
In our experiment, we used the IFBench dataset, which was divided into 40 development prompts and 260 test prompts. The experimental configuration and prompt hashes were saved before the formal run.
\subsection{Model Calls and Recovery}
The formal run used the gpt-4.1 and gpt-4o models. We set the temperature of 0, a random seed of 42, and an initial output limit of 4,096 tokens.

The experiment contained 2,080 expected identities: 260 test prompts, two models, and four conditions. The initial run produced 1,891 successful responses, 53 empty responses, and 136 errors. We subsequently reran selected empty or failed identities with higher output limits of 16,384 or 32,768 tokens. The final merged results contained 2,072 scorable responses and 8 errors and we used this merged set throughout the report.
\subsection{Metrics, Statistical Analysis, and Content Review}
Our primary metric was strict prompt-level accuracy: a prompt passed only if all its constraints were satisfied. We also reported loose prompt-level accuracy as a reference.

We also calculated the repair rate and regression rate:

\textit{Repair rate:} The number of questions that B0 failed but were revised and passed, divided by the total number of questions that B0 failed initially.

\textit{Regression rate:} The number of questions that B0 passed initially but failed after revision, divided by the total number of questions that B0 passed initially.

Furthermore, after the experiment was completed, we conducted an exact cluster sign-flip test to verify the paired differences. The absolute value of the sum of the differences was used as the statistic. Holm's multiple comparison correction covered six comparisons of relative B0 among the commonly corrected groups, and four exploratory comparisons that were separately corrected.

These tests assume that, when the null hypothesis is true, the labels of the compared conditions within each cluster can be interchanged. The complete test results can be found in Appendix~\ref{app:statistics}. It should be noted that the p-values reported and the bootstrap confidence intervals come from different statistical methods.

The content review covered 2,080 experimental identities across 260 test prompts, two models, and four conditions. Of these, 2,072 had responses that could be judged, while 8 had no available response. The initial labels were subsequently reviewed and confirmed by one author.

The content labels are categorized as ``completed'', ``partially completed'', ``uncompleted'', and ``unable to determine''.

\section{Results and Discussion}
\subsection{Main Results}
Table I illustrates the prompt-level accuracy of both models under the four experimental conditions. B2 received the highest observed strict pass rate for both models: 62.9\% for gpt-4.1 and 67.2\% for gpt-4o.

\begin{table}[!t]
\caption{Prompt-level accuracy of both models.}
\label{tab:main-results}
\centering
\begin{tabular}{lrrrr}
\toprule
& \multicolumn{2}{c}{gpt-4.1} & \multicolumn{2}{c}{gpt-4o} \\
\cmidrule(lr){2-3}\cmidrule(lr){4-5}
Condition & Strict & Loose & Strict & Loose \\
\midrule
B0 & 53.3\% & 59.8\% & 57.1\% & 62.9\% \\
B1 & 56.4\% & 62.9\% & 60.6\% & 65.3\% \\
B2 & \textbf{62.9\%} & \textbf{67.2\%} & \textbf{67.2\%} & \textbf{71.0\%} \\
B3 & 59.1\% & 64.5\% & 64.9\% & 69.1\% \\
\bottomrule
\end{tabular}
\end{table}
\subsection{Effect of Verifier Feedback}
Compared with B0, the strict pass rate of B2 largely increased by 9.65 percentage points for gpt-4.1 and 10.04 percentage points for gpt-4o. The differences remained statistically significant after Holm correction with adjusted ($p$)-values of 0.00042 and 0.00078, respectively.

B3 also increased paired strict accuracy relative to B0 by 5.79 percentage points for gpt-4.1 and 7.72 percentage points for gpt-4o. Both differences remained statistically significant after correction.

Although B2 had a higher observed strict pass rate than B3, the direct paired comparisons did not show statistically significant differences. B2$-$B3 was 3.86 percentage points for gpt-4.1 and 2.32 percentage points for gpt-4o, with adjusted ($p$)-values of 0.28661 and 0.40503, respectively. These results do not establish that the two conditions were equivalent or that B3 was inferior.
\subsection{Repairs and Regressions}
In the observed paired results, B2 repaired more initially failing responses and caused fewer regressions than B1.

For gpt-4.1, B1 repaired 27 of 121 initially failing responses and caused 19 of 138 initially passing responses to fail. B2 repaired 32 of 121 initially failing responses and caused 7 of 138 initially passing responses to fail.

For gpt-4o, B1 repaired 25 of 111 initially failing responses and caused 16 of 148 initially passing responses to fail. B2 repaired 36 of 111 initially failing responses and caused 10 of 148 initially passing responses to fail.

These results indicate that the B2 revision strategy was useful for both correcting failing responses and preserving passing responses in this experiment.
\subsection{Content Review}
We conducted a content review because passing the rule checks does not always mean that a response has fully completed the original task. Among the 2,080 experimental identities, 2,046 responses were marked as complete, 9 as partial, and 17 as incomplete. The remaining 8 could not be judged because the responses were empty or resulted from errors.

From the 26 responses labelled partial or incomplete, we found that 7 partial and 12 incomplete answers still passed the strict verifier whereas 818 complete responses did not pass the strict check. Therefore, these findings show that a strict rule alone cannot determine whether a response adequately completes the original task.

We also compared the content completion rates across the four conditions. For gpt-4.1, the rates were 99.2\% for B0, B1, and B2, and 98.5\% for B3. For gpt-4o, the corresponding rates were 98.5\%, 98.8\%, 98.5\%, and 98.1\%. These results illustrate that B2 achieved a high task completion rate while improving the strict constraint compliance relative to B0.

\section{Case Study}
\subsection{Successful Repair with Overall Verifier Feedback}
For item 173, the prompt asked whether pointwise convergence of polynomials on \([0,1]\) implies uniform convergence. It also required stop words to account for no more than 20\% of the response. The gpt-4.1 initial answer gave a mathematical counterexample but failed the strict constraint check:

\noindent\textbf{B0 --- FAIL}
\begin{quote}
No. Let \(B_n(x)=x^n\). For each \(x<1\), \(B_n(x)\to0\), while \(B_n(1)=1\). Thus the pointwise limit is \(f(x)=0\) for \(x<1\) and \(f(1)=1\). However, \(\sup_{x\in[0,1]}|B_n(x)-f(x)|=1\) for every \(n\), so convergence is not uniform.
\end{quote}

B2 received feedback that at least one output constraint had failed. Its revised answer passed the strict check:

\noindent\textbf{B2 --- PASS}
\begin{quote}
No. Choose \(B_n(x)=x^n\). Pointwise limit: \(f(x)=0\) for \(x<1\), and \(f(1)=1\). Yet \(\sup_{x\in[0,1]}|B_n(x)-f(x)|=1\) for every \(n\). Therefore, convergence is not uniform.
\end{quote}

Both answers used the same counterexample and reached the same conclusion. In this case, the revised answer satisfied the output constraint while retaining the main mathematical argument.
\subsection{Regression after Self-Revision}
For item 287, the prompt asked the model to copy the words between indices 10 and 12 from the following question:
\begin{quote}
how to find what value in a quadratic function lands on a certain point.
\end{quote}

The gpt-4.1 initial answer passed the strict check:

\noindent\textbf{B0 --- PASS}
\begin{quote}
on a certain
\end{quote}

After receiving the self-revision instruction, B1 returned the following answer, which failed the strict check:

\noindent\textbf{B1 --- FAIL}
\begin{quote}
lands on a
\end{quote}

The revised answer selected a span one word earlier than the accepted span. In this case, self-revision changed an initially passing answer into one that failed the output constraint, indicating a regression issue.

\section{Limitations}
First of all, B1 and B2/B3 prompts are different in a way that B2/B3 prompts were explicitly allowed the model to leave it unchanged whereas B1 did not receive the same conditional instruction. The observed difference between these conditions cannot therefore be attributed to verifier feedback alone.

Secondly, the results from B3 branch did not show whether per-constraint feedback is as effective as overall feedback. Although the comparisons between B2 and B3 did not show statistically significant differences, further experiments are still needed with repeated generations or more detailed feedback which is able to clarify the situation when per-constraint feedback is useful.

\section{Conclusion}
In this project, we evaluated whether verifier feedback could improve instruction following through a single revision step on IFBench. Compared with the initial answer B0, B2 increased paired strict prompt-level accuracy by 9.65 percentage points for gpt-4.1 and 10.04 percentage points for gpt-4o. B2 also had lower observed regression rates than B1. However, the comparisons between B2 and B3 did not show statistically significant differences.

The content review showed that passing the output constraints did not always mean that a response fully completed the original task. Future work could use matched prompts to isolate the effect of verifier feedback and investigate whether repeated revisions or more detailed feedback improve instruction following performance.

\section*{Author Contributions}
\noindent\textbf{Zhang Yanxiao (23222069)}
\begin{itemize}
  \item Experimental design
  \item Result analysis
  \item Experimental code development
  \item Report writing
  \item Presentation preparation
  \item Overall review and verification
\end{itemize}

\section*{GenAI Use Declaration}
ChatGPT and Grok were used to assist with experimental code development and to explore and design the experimental workflow. The report was written independently by the author. Generative AI was then used to correct grammatical errors and improve sentence structure.

\bibliographystyle{IEEEtran}
\bibliography{references}

\appendices
% Generated from the saved experiment summaries; see source paths in Appendix C.
\section{Statistical Comparisons and Robustness}
\label{app:statistics}
Tables~\ref{tab:strict-tests} and~\ref{tab:content-tests} report the post-hoc two-sided exact cluster sign-flip tests for strict accuracy and content completion, respectively. Model labels 4.1 and 4o refer to gpt-4.1 and gpt-4o. Each comparison uses 259 paired prompts per model. Identified related tasks form clusters; other prompts form individual clusters. The test assumes condition labels are exchangeable within clusters under the null. The difference $\Delta$ and its confidence interval (CI) are in percentage points (pp). For each outcome, Holm correction is applied separately to six B1/B2/B3 versus B0 comparisons and four exploratory comparisons of B2 versus B1/B3. The statistical tests and correction families were specified after the experiment. An interval of [0, 0] reflects the observed resamples and does not establish equivalence.


\begin{table}[!t]
\caption{Strict accuracy: paired differences and post-hoc tests.}
\label{tab:strict-tests}
\centering\scriptsize
\begin{tabular}{llrrrr}
\toprule
Model & Contrast & $\Delta$ & 95\% CI & $p$ & Holm $p$ \\
\midrule

4.1 & B1-B0 & +3.09 & [-1.93, 8.11] & 0.30200 & 0.42205 \\

4.1 & B2-B0 & +9.65 & [5.04, 14.01] & 0.00007 & 0.00042 \\

4.1 & B3-B0 & +5.79 & [1.51, 9.73] & 0.01067 & 0.03202 \\

4o & B1-B0 & +3.47 & [-1.16, 8.49] & 0.21102 & 0.42205 \\

4o & B2-B0 & +10.04 & [5.40, 15.18] & 0.00016 & 0.00078 \\

4o & B3-B0 & +7.72 & [3.47, 11.72] & 0.00082 & 0.00329 \\

4.1 & B2-B1 & +6.56 & [2.26, 11.54] & 0.00948 & 0.03790 \\

4.1 & B2-B3 & +3.86 & [-0.39, 8.37] & 0.14331 & 0.28661 \\

4o & B2-B1 & +6.56 & [1.93, 11.97] & 0.01862 & 0.05587 \\

4o & B2-B3 & +2.32 & [-2.30, 7.39] & 0.40503 & 0.40503 \\

\bottomrule
\end{tabular}
\end{table}

\begin{table}[!t]
\caption{Content completion: paired differences and post-hoc tests.}
\label{tab:content-tests}
\centering\scriptsize
\begin{tabular}{llrrrr}
\toprule
Model & Contrast & $\Delta$ & 95\% CI & $p$ & Holm $p$ \\
\midrule

4.1 & B1-B0 & +0.00 & [0.00, 0.00] & 1.00000 & 1.00000 \\

4.1 & B2-B0 & +0.00 & [0.00, 0.00] & 1.00000 & 1.00000 \\

4.1 & B3-B0 & -0.77 & [-1.93, 0.00] & 0.50000 & 1.00000 \\

4o & B1-B0 & +0.39 & [0.00, 1.17] & 1.00000 & 1.00000 \\

4o & B2-B0 & +0.00 & [-1.16, 1.15] & 1.00000 & 1.00000 \\

4o & B3-B0 & -0.39 & [-1.92, 0.78] & 1.00000 & 1.00000 \\

4.1 & B2-B1 & +0.00 & [0.00, 0.00] & 1.00000 & 1.00000 \\

4.1 & B2-B3 & +0.77 & [0.00, 1.95] & 0.50000 & 1.00000 \\

4o & B2-B1 & -0.39 & [-1.56, 0.78] & 1.00000 & 1.00000 \\

4o & B2-B3 & +0.39 & [-0.78, 1.92] & 1.00000 & 1.00000 \\

\bottomrule
\end{tabular}
\end{table}

Table~\ref{tab:robustness} compares paired point estimates under available-pair, complete-four-condition, and missing-as-failure treatments. The last is a sensitivity scenario rather than an observed failure label. The 95\% intervals reported in the source results are unadjusted percentile cluster-bootstrap intervals based on 1,000 resamples and seed 42; they are not Holm-adjusted and are calculated by a different procedure from the sign-flip $p$-values.
\begin{table}[!t]
\caption{Missing-data sensitivity: paired differences (pp).}
\label{tab:robustness}
\centering\footnotesize
\begin{tabular}{llrrr}
\toprule
Model & Contrast & Available & Complete & Missing fail \\
\midrule

4.1 & B1-B0 & +3.09 & +3.09 & +3.08 \\

4.1 & B2-B0 & +9.65 & +9.65 & +9.62 \\

4.1 & B3-B0 & +5.79 & +5.79 & +5.77 \\

4.1 & B2-B1 & +6.56 & +6.56 & +6.54 \\

4.1 & B2-B3 & +3.86 & +3.86 & +3.85 \\

4o & B1-B0 & +3.47 & +3.47 & +3.46 \\

4o & B2-B0 & +10.04 & +10.04 & +10.00 \\

4o & B3-B0 & +7.72 & +7.72 & +7.69 \\

4o & B2-B1 & +6.56 & +6.56 & +6.54 \\

4o & B2-B3 & +2.32 & +2.32 & +2.31 \\

\bottomrule
\end{tabular}
\end{table}

\section{Revision Prompts}
\label{app:prompts}
For B0, the model received the unmodified IFBench user prompt. Each revision branch received the original user prompt, the B0 assistant response, and a new user message containing the following components in order. Sections labeled as templates below describe the message construction; the labels themselves were not sent to the model.

\subsection{Shared Revision Instruction}

\begin{quote}
Please review your previous answer against the original user message. Keep the original task and every output constraint from that message.
\end{quote}

\subsection{B1: Self-Revision}
\begin{quote}
Check the original output constraints yourself. Do not assume they were already satisfied.
\end{quote}

\subsection{B2: Overall Verifier Feedback}
The following message was selected using the strict verifier outcome:
\begin{quote}
Verifier (strict): PASSED all output constraints.
\end{quote}
\begin{quote}
Verifier (strict): FAILED. At least one output constraint was not satisfied.
\end{quote}

\subsection{B3: Per-Constraint Verifier Feedback}
The header was followed by one numbered line per constraint, in the original order. Each line contained PASS or FAIL and the corresponding requirement from the original prompt.
\begin{quote}
Verifier (strict), one line per constraint, in listed order:
\end{quote}

\subsection{Passing Answer and Output Rule}
For B2 and B3, the following instruction was added only when B0 passed the strict verifier:
\begin{quote}
The previous answer already passed the output constraints. You may leave the answer unchanged.
\end{quote}
All branches ended with:
\begin{quote}
Return only the final answer. Do not describe the revision.
\end{quote}

\end{document}
